% STATUS: DRAFT
\chapter{Why algebras?}\label{ch:why_algebras}

\begin{physmotiv}
Textbook quantum mechanics says: a system is a Hilbert space $\HH$, observables are self-adjoint operators on $\HH$, states are unit vectors (or density matrices). For a finite number of degrees of freedom this is airtight. But three of the most interesting situations in theoretical physics---\emph{infinite systems}, \emph{quantum field theory in a region}, and \emph{quantum gravity}---quietly break one or more of these assumptions. This chapter shows how, using examples simple enough to compute on the back of an envelope. The cure in all three cases is the same: promote the \emph{algebra of observables} to the primary object, and let Hilbert spaces and states be derived from it.
\end{physmotiv}

\section{The textbook framework and what it hides}

Let us recall the framework with its hidden assumptions made explicit.

\begin{enumerate}[label=(\roman*)]
\item There is a \emph{single} Hilbert space $\HH$ for the system.
\item \emph{Every} reasonable state is a vector (or density matrix) in $\HH$.
\item \emph{Every} bounded function of the observables is again an observable, and $\Tr$ makes sense: $\langle A\rangle=\Tr(\rho A)$.
\item The positions and momenta $x_i,p_i$ are represented \emph{uniquely} (up to unitary equivalence) by the Heisenberg relations.
\item Observables associated with a subsystem $S$ live on a tensor factor: $\HH=\HH_S\otimes\HH_{\bar S}$, and the reduced density matrix $\rho_S=\Tr_{\bar S}\rho$ has an entropy $-\Tr\rho_S\log\rho_S$.
\item Physical observables can be associated with local regions of space.
\end{enumerate}

Statement (iv) is the Stone--von Neumann theorem (\cref{ch:weyl}) and is true for finitely many degrees of freedom. Statements (i)--(vi) \emph{all} fail in at least one of the three situations below.

\section{Crack 1: infinite systems and inequivalent representations}

Take an infinite chain of spin-$\tfrac12$ particles, one at each site $j\in\Z$. The \emph{local observables} are those that act non-trivially on finitely many sites, e.g.\ $\sigma^z_j$, $\sigma^x_j\sigma^x_{j+1}$. These obviously form an algebra (sums and products of local observables are local). What is the Hilbert space?

A natural guess is $\HH=\bigotimes_{j}\C^2$. But an infinite tensor product of Hilbert spaces is not well defined until one specifies a reference vector. Consider the following product states, with $\theta\in[0,\pi/2]$:
\begin{equation}
\ket{\Psi_\theta}=\bigotimes_{j\in\Z}\bigl(\cos\theta\ket{\uparrow}+\sin\theta\ket{\downarrow}\bigr).
\end{equation}
Truncate to $N$ sites. The overlap of two such states is
\begin{equation}
\braket{\Psi_\theta}{\Psi_{\theta'}}_N=\bigl(\cos\theta\cos\theta'+\sin\theta\sin\theta'\bigr)^N=\cos(\theta-\theta')^N\xrightarrow{N\to\infty}0\qquad(\theta\ne\theta').
\end{equation}
So these states are exactly orthogonal in the limit, and no vector in a ``universal'' Hilbert space can interpolate between them using local operations: any \emph{local} operator changes only finitely many factors and so cannot map $\ket{\Psi_\theta}$ to $\ket{\Psi_{\theta'}}$. Each $\Psi_\theta$ generates its own Hilbert space $\HH_\theta$ by acting with local operators, and these spaces carry different, \emph{unitarily inequivalent}, representations of the same algebra of local observables. Physically, $\Psi_\theta$ has magnetization $\langle\sigma^z_j\rangle=\cos 2\theta$, a macroscopic quantity that cannot be changed by any finite number of local operations. These are different phases, or superselection sectors.

\begin{keyidea}
In infinite systems, the algebra of local observables is unambiguous, but the Hilbert space is not: there are many inequivalent representations, corresponding to different macroscopic phases. A formalism that starts from a Hilbert space has already made a choice that cannot be made in advance.
\end{keyidea}

The same phenomenon appears for free bosons in infinite volume: the Fock representations at different temperatures and densities are inequivalent (\cref{ch:inequiv,ch:kms}), and the Stone--von Neumann theorem fails for infinitely many degrees of freedom (\cref{ch:weyl}).

\section{Crack 2: quantum field theory}

\subsection{Haag's theorem}

In the textbook treatment of an interacting QFT one writes the interacting field as a unitary rotation of the free field, $\phi_{\rm int}(x)=U^{-1}\phi_{\rm free}(x)U$, and uses the interaction picture. \emph{Haag's theorem} (1955) states, roughly, that if the free and interacting fields, both Poincar\'e covariant with a unique translation-invariant vacuum, are related by such a unitary, then the interacting field is itself free (its vacuum correlation functions coincide with the free ones). In other words, the interacting vacuum lives in a Hilbert space that is \emph{unitarily inequivalent} to the Fock space of the free theory. This is the same phenomenon as in the spin chain, now in a relativistic setting: the Fock space of Weyl operators is not the unique representation of the canonical commutation relations once there are infinitely many degrees of freedom. Perturbative QFT works in spite of this because it only ever computes finite-order matrix elements, never the unitary $U$ itself.

\subsection{Entanglement in a region has no density matrix}

Take a QFT in Minkowski space in its vacuum $\Omega$. Let $R$ be a region (say, the right Rindler wedge) and $\A(R)$ the algebra of observables localized in $R$. In a lattice regularization one would write $\rho_R=\Tr_{\bar R}\ket\Omega\!\bra\Omega$ and $S_R=-\Tr\rho_R\log\rho_R$. This entropy is divergent in the continuum limit: for a field theory in $d$ spacetime dimensions it is
\begin{equation}
S_R\;\sim\;c_1\frac{\mathrm{Area}(\partial R)}{\epsilon^{d-2}}+\dots
\end{equation}
where $\epsilon$ is a short-distance cutoff (and, for a two-dimensional CFT, $S=\tfrac{c}{3}\log(L/\epsilon)$ for an interval of length $L$). The coefficient of the leading divergence is non-universal, it depends on how one cuts off the theory. The divergence is a symptom of something deeper than a UV nuisance:

\begin{itemize}
\item By the \emph{Reeh--Schlieder theorem} (\cref{ch:rs_bw}) the vacuum $\Omega$ is \emph{cyclic and separating} for $\A(R)$: acting on $\Omega$ with local operators in $R$ one can approximate any state in $\HH$, and no non-zero operator in $\A(R)$ annihilates $\Omega$. This means the vacuum is very entangled between $R$ and its complement, to the extent that no ``reduced density matrix'' exists.
\item The algebra $\A(R)$ is not of the form $\BH_R\otimes\id$ for any factor $\HH_R$. It is a \emph{type III$_1$} von Neumann algebra (\cref{ch:typeIII_local}), which has no trace, no pure states, and no density matrices.
\end{itemize}

So the question ``what is the entropy of a region in QFT?'' has no canonical answer in the continuum, and the reason is algebraic: the algebra of a region is of a type that does not admit entropy. One of the aims of this book is to see exactly how gravity changes this (\cref{ch:crossed,ch:entropy_II,ch:clpw}).

\section{Crack 3: gravity}

In a diffeomorphism-invariant theory, gauge invariance is the statement that physical observables commute with the constraints that generate diffeomorphisms. A local field $\phi(x)$ at a coordinate point $x$ is not gauge invariant: moving the point is a gauge transformation. Only \emph{relational} observables (``the value of $\phi$ where the curvature scalar $=1$'', ``the value of $\phi$ one proper time $\tau$ along the worldline of the observer'') are meaningful.

Two consequences matter for us.

\begin{enumerate}
\item The algebra of observables of a spatial region is no longer simply the algebra generated by the local fields in that region. One must ask: relative to \emph{what} is the region defined? In a closed universe, such as de Sitter space with its compact spatial sections, there is no asymptotic boundary relative to which to define operators.
\item Imposing the constraints changes the \emph{type} of algebra. We will see that the algebra of observables of a static patch in de Sitter space, once an observer is included, is a type II$_1$ factor. Unlike type III$_1$, this algebra has a trace and hence a well-defined entropy, and it has a maximal-entropy state: the de Sitter vacuum.
\end{enumerate}

A precursor of this phenomenon is already present in ordinary gauge theory, where the constraint (Gauss law) forces certain edge degrees of freedom to be associated to a region's boundary and changes the factorization of the Hilbert space (\cref{ch:gauge_warmup}).

\section{The algebraic viewpoint}

The common remedy is to take the following dictionary seriously.

\begin{center}
\renewcommand{\arraystretch}{1.3}
\begin{tabular}{@{}ll@{}}
\toprule
Textbook QM & Algebraic formulation\\
\midrule
Observables = self-adjoint operators on $\HH$ & Observables = self-adjoint elements of a $^*$-algebra $\A$\\
State = unit vector / density matrix & State = positive normalized linear functional $\omega:\A\to\C$\\
Hilbert space $\HH$ (given) & Hilbert space $\HH_\omega$ (constructed from $\omega$ by GNS, \cref{ch:gns})\\
Unique $x,p$ representation & Many inequivalent representations\\
Subsystem = tensor factor & Subsystem = subalgebra\\
Locality = tensor factorization & Locality = commutation of algebras of spacelike regions\\
\bottomrule
\end{tabular}
\end{center}

The expectation value $\omega(A)$ of an observable $A\in\A$ in a state $\omega$ is the only input from experiment. For a finite system $\A=\BH$ and $\omega(A)=\Tr(\rho A)$, so nothing is lost. For an infinite system the dictionary is the correct generalization, as the next chapters show.

The algebraic viewpoint also has a practical benefit that physicists will recognize: it separates the \emph{kinematics} (the algebra, determined by locality, symmetry and commutation relations) from the \emph{dynamics} and \emph{state}. Many properties of QFT, such as the type of the local algebras, do not depend on the details of the interaction at all.

\begin{whatwrong}
``Algebraic'' does not mean ``Hilbert-space free''. In practice we always represent $\A$ on a Hilbert space, and a vN algebra is by definition an algebra of operators on one. The point is that no particular Hilbert space is privileged, and that the properties of the algebra that matter (e.g.\ its type) are independent of which faithful representation is chosen.
\end{whatwrong}

\section{What this book assumes and does not assume}

We assume graduate-level quantum mechanics, QFT (canonical quantization, Wick rotation, Rindler space), some statistical mechanics (Gibbs states, partition functions), and general relativity at the level of Schwarzschild black holes and horizons. We assume \emph{no} functional analysis or topology. Every mathematical concept is introduced after a physical motivation; we prove what is illuminating, sketch what is not, and cite what is heavy.

\section*{Exercises}
\begin{exercise}
Verify the overlap formula $\braket{\Psi_\theta}{\Psi_{\theta'}}_N=\cos(\theta-\theta')^N$. Show that the expectation value of any local operator in $\Psi_\theta$ is well defined (depends on a finite number of sites). Compute the magnetization $\langle\sigma^z_j\rangle$.
\end{exercise}
\begin{exercise}
Let $\ket{\psi}\in\C^2\otimes\C^2$ be $\cos\alpha\ket{\uparrow\uparrow}+\sin\alpha\ket{\downarrow\downarrow}$. Compute $\rho_A$ and its entropy. How does the entropy behave as $\alpha\to0$? Now consider $N$ copies of this state. In what sense does a limit $N\to\infty$ of the entanglement entropy fail to exist (for fixed $\alpha\ne0$)?
\end{exercise}
\begin{exercise}
For a free massless scalar on a lattice in $1+1$ dimensions with spacing $\epsilon$, the vacuum entanglement entropy of an interval of $L/\epsilon$ sites scales as $\tfrac13\log(L/\epsilon)$. Explain why the continuum limit $\epsilon\to0$ at fixed $L$ cannot exist, and why this does not contradict the existence of well-defined correlation functions in the continuum.
\end{exercise}
\begin{exercise}
Argue (without computation) that gauge-invariant observables in a theory with a Gauss law cannot be written as a tensor product over the two sides of a surface. (We return to this in \cref{ch:gauge_warmup}.)
\end{exercise}
